\documentclass[12pt,a4paper]{article}
\usepackage[utf8]{inputenc}
\usepackage[T1]{fontenc}
\usepackage[brazil]{babel}
\usepackage{amsmath,amssymb}
\usepackage[margin=2.5cm]{geometry}

\begin{document}

\section{Definição da Probabilidade $p$}
Para atribuir uma probabilidade específica a equipe, considera-se o
maior número de matrícula entre os integrantes do grupo. Os números de
matrícula são apresentados na Tabela~\ref{tab:matriculas}.

\begin{table}[h]
\centering
\begin{tabular}{lc}
\hline
\textbf{Integrante} & \textbf{Matrícula} \\
\hline
Allan Martins Gadelha & 579087 \\
Clarisse Maria Cabral Montenegro & 580585 \\
Maria Luisa Fernandes de Mendonça & 581706 \\
Sara de Sousa Magalhães & 581189 \\
\hline
\end{tabular}
\caption{Números de matrícula dos integrantes do grupo.}
\label{tab:matriculas}
\end{table}

A maior matrícula entre os integrantes é:
\[
M = 581706.
\]
Conforme a definição estabelecida, calcula-se $p$ por meio da expressão:
\[
p = 0{,}05 + \frac{M \bmod 100}{200}.
\]
O operador $\bmod$ representa o resto da divisão inteira de $M$ por 100.
Calculando, obtém-se:
\[
581706 \bmod 100 = 6.
\]
Logo,
\[
p = 0{,}05 + \frac{6}{200} = 0{,}05 + 0{,}03 = \mathbf{0{,}08}.
\]


\end{document}
